\chapter{Spring Kafka: Producer, Listener, Poison Pill}
\label{ch:springkafka}

\section{Phát sự kiện}

Phát một sự kiện trong dự án này gồm hai bộ phận: một
\code{KafkaTemplate} và cấu hình chọn serializer. Từ auth-service:

\begin{lstlisting}[language=Java]
kafkaTemplate.send("ts.user.registered", event.userId(), event);
//                  topic                key             value
\end{lstlisting}

\begin{lstlisting}[language=yaml]
spring:
  kafka:
    bootstrap-servers: localhost:9092
    producer:
      key-serializer: org.apache.kafka.common.serialization.StringSerializer
      value-serializer: org.springframework.kafka.support.serializer.JsonSerializer
\end{lstlisting}

Key (\code{userId}) là một chuỗi; value là record sự kiện được tuần tự
hóa thành JSON. Dùng user id làm key là quyết định về thứ tự ở Chương~10
được cụ thể hóa. \code{send()} là bất đồng bộ --- nó trả về một future;
thông điệp nằm trong bộ đệm phía client cho tới khi broker ack. Hai hệ
quả vận hành: một service có thể \emph{khởi động} mà không có broker
(các lần gửi thất bại về sau, theo từng lời gọi --- đó là lý do
auth-service có thể triển khai lên cụm trước Kafka), và một service sập
ngay sau khi \code{send()} trả về có thể mất thông điệp đang nằm trong
bộ đệm --- mẫu outbox bên dưới tồn tại cho khi điều đó quan trọng.

\section{Tiêu thụ}

\begin{lstlisting}[language=Java]
@KafkaListener(topics = "ts.user.registered",
               groupId = "notification-group")
public void on(UserRegisteredEvent event) {
    mailService.sendActivation(event);
}
\end{lstlisting}

Spring chạy một vòng lặp poll cho mỗi listener container, giải tuần tự
từng record, gọi phương thức, và commit offset sau khi xử lý thành công
--- ít nhất một lần ngay từ đầu. Nửa consumer của file YAML chứa những
dòng liên quan tới production nhất của chương này:

\begin{lstlisting}[language=yaml]
spring:
  kafka:
    consumer:
      group-id: course-group
      auto-offset-reset: earliest                    (*@\co{1}@*)
      key-deserializer: ...ErrorHandlingDeserializer (*@\co{2}@*)
      value-deserializer: ...ErrorHandlingDeserializer
      properties:
        spring.deserializer.value.delegate.class: ...JsonDeserializer
        spring.json.trusted.packages: com.ts.common.dto  (*@\co{3}@*)
\end{lstlisting}

\begin{description}
  \item[\co{1}] Nơi một group \emph{hoàn toàn mới} bắt đầu khi chưa có
  offset đã commit: \code{earliest} = từ đầu khoảng retention; mặc định
  \code{latest} = chỉ thông điệp mới. \code{earliest} nghĩa là một
  consumer triển khai muộn vẫn xử lý các sự kiện đã phát trước khi nó tồn
  tại --- lựa chọn đúng khi sự kiện là những sự thật cần hành động. Không
  còn ý nghĩa một khi group đã có offset.
  \item[\co{2}] Lá chắn chống poison pill --- mục tiếp theo.
  \item[\co{3}] Giải tuần tự JSON khởi tạo các lớp được nêu tên trong
  header thông điệp; chỉ tin cậy \code{com.ts.common.dto} ngăn một
  producer độc hại hoặc có lỗi khiến consumer khởi tạo lớp tùy ý. Một
  ranh giới bảo mật, không phải mã lặp.
\end{description}

\section{Poison pill}

Một \emph{poison pill} (viên thuốc độc) là thông điệp mãi mãi không giải
tuần tự được --- lỗi ở producer, một trường bị đổi tên, payload hỏng.
Không có phòng thủ, listener ném lỗi \emph{trước khi} thông điệp có thể
bị bỏ qua, offset không bao giờ tiến lên, và container thử lại cùng một
record không ngừng: partition bị kẹt và log đầy lên với tốc độ tối đa.

\code{ErrorHandlingDeserializer} bọc deserializer thật: khi thất bại nó
phát ra một lỗi có kiểu thay vì ném ngoại lệ bên trong vòng lặp poll, để
error handler của Spring bỏ qua record đó (tùy chọn sau vài lần thử lại,
tùy chọn phát nó sang một dead-letter topic). Cách ghép --- error handler
bên ngoài, JSON delegate bên trong --- là mẫu production chuẩn, và là lý
do cấu hình nêu tên \emph{hai} lớp deserializer cho mỗi bên.

\begin{pitfall}[một thông điệp hỏng, im lặng hoàn toàn]
Triệu chứng: email thông báo dừng hẳn; consumer ghi cùng một stack trace
giải tuần tự hàng nghìn lần. Nguyên nhân: một poison pill trên một
partition mà không có deserializer xử lý lỗi --- vòng lặp không xử lý
được cũng không bỏ qua được. Lỗi này ẩn mình trong những hệ thống chưa
bao giờ thử với một thông điệp sai định dạng. Hãy kiểm chứng lá chắn của
bạn một lần, có chủ đích: phát rác vào một topic dev và xem nó bị bỏ
qua.
\end{pitfall}

\section{Những góc nâng cao}

\begin{description}
  \item[Dead-letter topic.] \code{DefaultErrorHandler} +
  \code{DeadLetterPublishingRecordRecoverer}: sau N lần thử thất bại,
  record rơi vào \code{<topic>.DLT} để xem xét và phát lại --- xử lý lỗi
  có dấu vết kiểm toán.
  \item[Mẫu outbox.] ``Lưu người dùng \emph{và} phát sự kiện'' là hai hệ
  thống và không thể nguyên tử một cách trực tiếp. Outbox: ghi sự kiện
  vào một bảng \code{outbox} trong \emph{cùng transaction cơ sở dữ liệu}
  với hàng người dùng; một relay (hoặc Debezium đọc WAL) phát từ bảng đó.
  Nếu một ngày auth-service không được phép mất sự kiện đăng ký, đây là
  câu trả lời.
  \item[Schema registry.] Các lớp sự kiện dùng chung ở Chương~6 giải
  quyết việc producer/consumer thống nhất với nhau bên trong một repo. Ở
  quy mô tổ chức, hợp đồng chuyển vào một schema registry (schema
  Avro/Protobuf, được kiểm tra tương thích lúc phát) --- cùng ý tưởng,
  nhưng do hạ tầng áp đặt thay vì một jar dùng chung.
\end{description}

\begin{quiz}
\begin{enumerate}
  \item Vì sao auth-service khởi động được khi Kafka đang sập, và đến
  thời điểm nào thì một sự cố Kafka trở nên hữu hình với nó?
  \item Một consumer group mới được triển khai hôm nay. Giải thích
  \code{auto-offset-reset: earliest} khiến nó làm gì, và khi nào thiết
  lập này không còn tác dụng với group đó.
  \item Đi qua kịch bản poison pill khi không có
  \code{ErrorHandlingDeserializer}, rồi khi có. Trong một hệ thống
  production đầy đủ, record hỏng cuối cùng nằm ở đâu?
  \item Vì sao \code{spring.json.trusted.packages} là một biện pháp kiểm
  soát bảo mật?
  \item Khi nào mẫu outbox là cần thiết, và nó làm cho cặp thao tác ghi
  nào trở nên nguyên tử một cách hiệu quả?
\end{enumerate}
\end{quiz}
